\documentclass[10pt,a4paper]{article}
\usepackage[top=1.5cm, bottom=1.5cm, left=1.8cm, right=1.8cm]{geometry}
\usepackage{booktabs}
\usepackage{array}
\usepackage{xcolor}
\usepackage{enumitem}
\usepackage{parskip}
\usepackage[T1]{fontenc}
\usepackage{microtype}

\setlength{\parskip}{3pt}
\pagestyle{empty}

\begin{document}

\begin{center}
{\large\textbf{COLDSTART: Cost- and Latency-Optimized Serverless Orchestration}}\\[2pt]
{\normalsize\textbf{via Multi-Modal Prediction and Confidence-Based Warming}}\\[4pt]
\small
Abhishek Doke \quad Nikhil Dhamdhere \quad Joanna Dharwadkar \quad Geetanjali Bejjalwar\\
\textit{Guide: Dr. Alpana Adsul}\\
Department of Computer Engineering, Dr. D.Y. Patil College of Engineering and Innovation, Varale, Talegaon, Pune \quad 2024--2025
\end{center}

\noindent\rule{\linewidth}{0.8pt}

\smallskip
\noindent\textbf{Problem Statement:}
AWS Lambda (Function-as-a-Service) suffers from \textit{cold start latency} of 500\,ms to several seconds when an idle function receives a new request. Existing solutions use static keep-alive policies or single-model predictions, causing 40--60\% resource waste. There is no system that simultaneously minimizes cold start latency \textit{and} resource cost using confidence-aware ML predictions.

\smallskip
\noindent\textbf{Proposed Solution:}
COLDSTART is a cloud-native intelligent orchestration system that proactively predicts cold start events using a dual-model ensemble (XGBoost + Random Forest) and automatically pre-warms Lambda functions before requests arrive. A novel \textbf{Confidence-Based Warming (CBW)} mechanism gates pre-warming only when both models agree with high confidence, eliminating wasteful resource allocation.

\smallskip
\noindent\textbf{System Components:}
\begin{itemize}[leftmargin=1.4em, itemsep=1pt, topsep=1pt]
\item \textbf{Dual-Model Ensemble:} XGBoost (70\%) + Random Forest (30\%) on AWS SageMaker Serverless Endpoints predict cold start probability in real time from CloudWatch metrics.
\item \textbf{IntelligentLambdaWarmer:} AWS Lambda function that evaluates predictions, applies CBW gating, and issues pre-warm pings to target functions (auth, profile, search).
\item \textbf{ETL Pipeline:} Automated DynamoDB-to-S3 pipeline processes traffic data every 15 minutes to generate ML input features.
\item \textbf{Monitoring Dashboard:} Streamlit-based real-time dashboard displaying cold start probability, warming success rate, cost savings, and system health via CloudWatch API.
\item \textbf{Traffic Generator:} EC2-based load simulator generating realistic API traffic patterns for system testing and validation.
\end{itemize}

\smallskip
\noindent\textbf{Technology Stack:}

\noindent
\begin{tabular}{@{} p{3cm} p{12.5cm} @{}}
Cloud Platform & AWS Lambda, SageMaker, CloudWatch, DynamoDB, S3, EC2, API Gateway \\
ML / Backend   & Python 3.x, Scikit-learn, XGBoost, Random Forest, Boto3, Pandas \\
Frontend       & Streamlit Dashboard, AWS CloudWatch API \\
\end{tabular}

\smallskip
\noindent\textbf{Key Results:}

\noindent
\begin{tabular}{@{} p{5cm} p{2.2cm} p{5cm} p{2.2cm} @{}}
\toprule
\textbf{Metric} & \textbf{Value} & \textbf{Metric} & \textbf{Value} \\
\midrule
Cold Start Reduction   & 58.7\%  & Prediction Accuracy (F1) & 94.2\% \\
P95 Latency            & 187\,ms & Warming Success Rate      & 98.1\% \\
Cost Efficiency        & 34.8\%  & System Uptime             & 99.3\% \\
\bottomrule
\end{tabular}

\smallskip
\noindent\textbf{Novelty / Original Contribution:}
The Confidence-Based Warming (CBW) mechanism is the core original contribution. It calculates a confidence score from the agreement between XGBoost and Random Forest predictions. Pre-warming is triggered only when both the cold start probability exceeds a threshold \textit{and} the model confidence is high, directly minimizing false positives (wasted warming) while maintaining low latency. This dual-condition gating policy is not present in any existing cold start mitigation system. Unlike prior work that uses static keep-alive timers or single-model predictions, COLDSTART dynamically adapts its warming decisions based on real-time ensemble consensus, making it both cost-aware and performance-aware simultaneously.

\smallskip
\noindent\textbf{Comparison with Existing Approaches:}

\noindent
\begin{tabular}{@{} p{4.2cm} p{2.8cm} p{2.8cm} p{2.8cm} p{2.8cm} @{}}
\toprule
\textbf{Feature} & \textbf{Static Keep-Alive} & \textbf{Single ML Model} & \textbf{Priority Heuristic} & \textbf{COLDSTART} \\
\midrule
Cold Start Reduction  & Low     & Moderate & High     & \textbf{Highest} \\
Resource Waste        & High    & High     & Moderate & \textbf{Lowest}  \\
Confidence Gating     & No      & No       & No       & \textbf{Yes}     \\
Real-time Dashboard   & No      & No       & No       & \textbf{Yes}     \\
AWS Native Deployment & Partial & Partial  & No       & \textbf{Full}    \\
\bottomrule
\end{tabular}

\smallskip
\noindent\textbf{Future Scope:}
\begin{itemize}[leftmargin=1.4em, itemsep=1pt, topsep=1pt]
\item Extend to multi-cloud platforms (Azure Functions, Google Cloud Run) with a unified telemetry abstraction layer.
\item Replace static ensemble weights with a Deep Reinforcement Learning agent that continuously optimizes warming thresholds based on live cost-benefit feedback.
\item Add container memory and CPU size prediction to move beyond binary warm/no-warm decisions toward optimal resource sizing.
\item Integrate anomaly detection to identify unusual traffic spikes and trigger emergency pre-warming before SLA violations occur.
\end{itemize}

\noindent\rule{\linewidth}{0.8pt}
\begin{center}
\small\textit{Final Year Project --- B.E. Computer Engineering --- Academic Year 2024--2025}
\end{center}

\end{document}
